\section{Advantages}

\subsection{Distribution Independence}

The gPCD method is effective for dispersed, sparse, dense, static, and fast-moving particle populations. Because occupancy assignment is determined directly from the instantaneous particle geometry, the method does not rely on temporal coherence, cached neighbor lists, or motion-history assumptions. Consequently, performance remains stable when particles are widely separated, rapidly translating, or undergoing frequent redistribution.

This allows scarce, mobile, and stationary particle systems to be processed with comparable efficiency, provided local cell occupancy remains bounded.

\subsection{Constant-Time Cell Access}

Memory access to an individual cell remains constant-time,

\[
\mathcal{O}(1),
\]

because each cell is addressed through deterministic index mapping rather than through iterative search or hierarchical traversal. Cell retrieval therefore remains predictable and independent of the total number of cells.

Although the total storage capacity of the cell array may scale as

\[
\mathcal{O}(N^3),
\]

for a three-dimensional domain discretization, this increase in global memory allocation does not affect per-access lookup cost.

\subsection{Broad-Phase and Graphics Concurrency}

Broad-phase occupancy generation is executed concurrently with graphics-pipeline activity through the vertex and fragment stages. As a result, collision candidate generation is overlapped with normal rendering operations rather than requiring an entirely separate preprocessing pass.

This hybrid use of rasterization and compute resources improves hardware utilization and reduces the effective standalone overhead of broad-phase processing.

\subsection{Bounded Neighbor Search}

Under the particle-size constraint,

\[
d_p \leq 1,
\]

each particle may intersect no more than eight neighboring cells. This imposes a deterministic upper bound on per-particle cell references and creates predictable memory-access behavior during narrow-phase processing.

\subsection{Duplicate-Resistant Candidate Evaluation}

The method removes both corner duplicates and repeated particle-pair duplicates through local per-thread duplicate filtering. This reduces redundant collision checks without requiring global synchronization, improving narrow-phase efficiency and preserving parallel scalability.



